%% LyX 2.3.1 created this file.  For more info, see http://www.lyx.org/.
%% Do not edit unless you really know what you are doing.
\documentclass[english,final]{beamer}
\usepackage{fourier}
\renewcommand{\familydefault}{\rmdefault}
\usepackage[T1]{fontenc}
\usepackage[latin9]{inputenc}
\setcounter{secnumdepth}{3}
\setcounter{tocdepth}{3}
\setlength{\parskip}{\medskipamount}
\setlength{\parindent}{0pt}
\usepackage{url}
\usepackage{multicol}
\usepackage{graphicx}

\makeatletter

%%%%%%%%%%%%%%%%%%%%%%%%%%%%%% LyX specific LaTeX commands.


\title[Template]{Example Beamer Poster}


\subtitle{One Giant Slide, Subdivided Various Ways}


\author[PJ]{Paul E. Johnson}


\institute[CRMDA]{CRMDA, University of Kansas}


\date{2018}


\titlegraphic{\medskip{}
\includegraphics[width=16cm]{static/logo}}


\newcommand{\noun}[1]{\textsc{#1}}
\DeclareRobustCommand*{\lyxarrow}{%
\@ifstar
{\leavevmode\,$\triangleleft$\,\allowbreak}
{\leavevmode\,$\triangleright$\,\allowbreak}}
%% Because html converters don't know tabularnewline
\providecommand{\tabularnewline}{\\}

%%%%%%%%%%%%%%%%%%%%%%%%%%%%%% Textclass specific LaTeX commands.
\usepackage{etoolbox}
\AtEndPreamble{%
   \usepackage{beamerposter}
}
\providecommand*{\code}[1]{\texttt{#1}}
\providecommand*{\strong}[1]{\textbf{#1}}

%%%%%%%%%%%%%%%%%%%%%%%%%%%%%% User specified LaTeX commands.
\usepackage{xcolor}

%KU color pallette
%primary
\definecolor{kublue}{RGB}{0,81,186}
\definecolor{kucrimson}{RGB}{232,0,13}
\definecolor{kujayhawkyellow}{RGB}{255,200,45}
\definecolor{kusignaturegrey}{RGB}{133,137,138}
%secondary
\definecolor{kunight}{RGB}{0,52,89}
\definecolor{kulake}{RGB}{39,103,255}
\definecolor{kusky}{RGB}{115,203,242}
\definecolor{kubrick}{RGB}{151,27,47}
\definecolor{kufire}{RGB}{255,48,66}
\definecolor{kuwheat}{RGB}{242,169,0}
\definecolor{kufog}{RGB}{142,159,188}
\definecolor{kusteam}{RGB}{221,229,237}
\definecolor{kuterracotta}{RGB}{192,110,78}
\definecolor{kulimestone}{RGB}{215,210,203}


% You can select your theme here. Any beamer theme without a sidebar should work:
%\usetheme{Berlin}
\useoutertheme[footline=authorinstitutetitle]{miniframes}
\useinnertheme{rounded}
\setbeamercolor{structure}{fg=kublue}
%\usecolortheme{whale}
%\usecolortheme{orchid}
%replace some colors

\setbeamercolor{frametitle}{bg=kublue,fg=white}
\setbeamercolor{structure}{bg=kublue,fg=kublue!90!black}
%disable rounded for alert and example boxes%
\setbeamertemplate{blocks}[default]
\setbeamercolor*{palette primary}{use=structure,fg=white,bg=structure.fg}
\setbeamercolor*{palette secondary}{use=structure,fg=white,bg=kulake}
\setbeamercolor*{palette tertiary}{use=structure,fg=white,bg=kublue}
\setbeamercolor*{palette quaternary}{fg=white,bg=black}

\setbeamercolor*{sidebar}{use=structure,bg=structure.fg}
\setbeamercolor*{palette sidebar primary}{use=structure,fg=structure.fg!10}
\setbeamercolor*{palette sidebar secondary}{fg=white}
\setbeamercolor*{palette sidebar tertiary}{use=structure,fg=structure.fg!50}
\setbeamercolor*{palette sidebar quaternary}{fg=white}

\setbeamercolor*{titlelike}{parent=palette primary}

\setbeamercolor*{separation line}{}
\setbeamercolor*{fine separation line}{}

\setbeamercolor{block title}{use=structure,fg=white,bg=structure.fg}
\setbeamercolor{block title alerted}{use=alerted text,fg=white,bg=alerted text.fg!75!black}
\setbeamercolor{block title example}{use=example text,fg=white,bg=example text.fg!75!black}

\setbeamercolor{block body}{parent=normal text,use=block title,bg=block title.bg!10!bg}
\setbeamercolor{block body alerted}{parent=normal text,use=block title alerted,bg=block title alerted.bg!10!bg}
\setbeamercolor{block body example}{parent=normal text,use=block title example,bg=block title example.bg!10!bg}



% Note, however, that only dedicated beamerposter themes also generate
% a suitable poster title automatically (we do this manually, see the notes
% in the document). 
% Dedicated beamerposter themes are available at 
% http://www-i6.informatik.rwth-aachen.de/~dreuw/latexbeamerposter.php

% Here you can set the options (note that the example content is made for the given setup. 
% With different orientation or poster size, the content probably won't fit on the poster):
\PassOptionsToPackage{%
   orientation=portrait, % possible: portrait, landscape
    size=a0,% possible: a0b, a0, a1, a2, a3, a4, custom (with width=<val>,height=<val>)
    scale=1.4,% scaling of fonts
    debug% give verbose warnings and error messages
}{beamerposter}

% This produces a grid which might be useful when positioning things
%\beamertemplategridbackground[0.5cm]
%\setbeamertemplate{background canvas}[vertical shading][bottom=gray,top=white]
\usepackage{tcolorbox}


\definecolor{light-gray}{gray}{0.90}
\usepackage{realboxes}
\providecommand*{\code}[1]{\texttt{#1}}
\renewcommand{\code}[1]{%
%\Colorbox{light-gray}{#1}%
\Colorbox{kulimestone}{#1}%
}%

\makeatother

\usepackage{babel}
\begin{document}
\begin{frame}

\parbox[t]{0.8\columnwidth}{%
% Output title
{\VERYHuge\textbf{\inserttitle}}
\ifx\insertsubtitle\empty\else%
\par\vspace{1.2cm}{\veryHuge\insertsubtitle}\fi\vspace{1.5cm}

% Output author and date
{\veryHuge\insertauthor\ $\cdot$ \insertdate}%
}%
\parbox[t]{0.2\columnwidth}{%
% Logo
\raggedleft\inserttitlegraphic%
}

\vspace{1.7cm}

\begin{block}{One Block that is as wide as the whole poster (probably wasting space)}

This commentary is inside the larger blue block.There is also a Green,
Red, and Blue box within this one super-wide block.
\begin{exampleblock}{Example blocks appear as green}

The only thing that is different about the Example Block is that its
outline is GREEN. Equations work as expected, but cause a lot of empty
space within this giant wide block(you could check that).

\end{exampleblock}
\begin{alertblock}{An Alert Block is Red}

The red text draws attention. 

\end{alertblock}
\begin{block}{This is a standard blue block}

This is a standard blue block. After this, we create space with a
vertical fill. Everything from ``One Block that is as wide as the
whole poster'' to here is inside the wide block.
\end{block}
\end{block}
\vspace{1cm}

\begin{block}{A Four Column Area within a Block}

\setlength{\columnsep}{2.2cm}

\begin{multicols}{4}

Wide blocks are difficult to read. One way to solve this is using
the LaTeX \code{\textbackslash multicols} environment. LyX has a
module, ``Multiple Columns'' to facilitate (documented in \structure{Help\lyxarrow Additional Features}).
The content will automatically ``snake'' from top left to fill the
columns. The space between the columns can be adjusted via the \LaTeX{}
code \alert{\textbackslash setlength\{\textbackslash columnsep\}}
. A manual column break is inserted here with \code{\textbackslash vfill\textbackslash null}
followed by a \alert{\textbackslash columnbreak}.\vfill\null

\columnbreak

Lorem ipsum dolor sit amet, consetetur sadipscing elitr, sed diam
nonumy eirmod tempor invidunt ut labore et dolore magna aliquyam erat,
sed diam voluptua. At vero eos et accusam et justo duo dolores et
ea rebum. Stet clita kasd gubergren, no sea takimata sanctus est Lorem
ipsum dolor sit amet. Lorem ipsum dolor sit amet, consetetur sadipscing
elitr, sed diam nonumy eirmod tempor invidunt ut labore et dolore
magna aliquyam erat, sed diam voluptua. At vero eos et accusam et
justo duo dolores et ea rebum. Stet clita kasd gubergren, no sea takimata
sanctus est Lorem ipsum dolor sit amet. Lorem ipsum dolor sit amet,
consetetur sadipscing elitr, sed diam nonumy eirmod tempor invidunt
ut labore et dolore magna aliquyam erat, sed diam voluptua. At vero
eos et accusam et justo duo dolores et ea rebum. Stet clita kasd gubergren,
no sea takimata sanctus est Lorem ipsum dolor sit amet.

Duis autem vel eum iriure dolor in hendrerit in vulputate velit esse
molestie consequat, vel illum dolore eu feugiat nulla facilisis at
vero eros et accumsan et iusto odio dignissim qui blandit praesent
luptatum zzril delenit augue duis dolore te feugait nulla facilisi.
Lorem ipsum dolor sit amet, consectetuer adipiscing elit, sed diam
nonummy nibh euismod tincidunt ut laoreet dolore magna aliquam erat
volutpat.

Ut wisi enim ad minim veniam, quis nostrud exerci tation ullamcorper
suscipit. 

\begin{tcolorbox}[title=Heres is a tcolorbox object that seems to fig gracefully within a column]

\setlength{\columnsep}{2.2cm}

I'm shocked this works.  When I put a beamer block in here, it ran off into the right margin!.
\end{tcolorbox}
\end{multicols}

\end{block}

\begin{columns}[t, totalwidth=1.0105\textwidth]

\column{.49\linewidth}
\begin{block}{Fontsizes: From tiny to most giant!}

\medskip{}

\begin{tabular}{cc}
{\tiny{}Tiny }(\texttt{\textbackslash tiny}) & {\scriptsize{}Smallest }(\texttt{\textbackslash scriptsize})\tabularnewline
{\footnotesize{}Smaller }(\texttt{\textbackslash footnotesize}) & Normal (\texttt{\textbackslash normalsize})\tabularnewline
{\large{}Large }(\texttt{\textbackslash large}) & {\Large{}Larger }(\texttt{\textbackslash Large})\tabularnewline
{\LARGE{}Largest }(\texttt{\textbackslash LARGE}) & {\huge{}Huge }(\texttt{\textbackslash huge})\tabularnewline
 & \multicolumn{1}{c}{}\tabularnewline
{\Huge{}Huger }(\texttt{\textbackslash Huge}) & {\veryHuge{}Giant \par}(\texttt{\textbackslash veryHuge})\tabularnewline
\multicolumn{2}{c}{{\VERYHuge{}Gigantic \par}(\textbackslash VERYHuge)}\tabularnewline
\end{tabular}

\end{block}
\begin{alertblock}{Blocks in different colors}
\begin{alertblock}{}
Beamer provides special blocks with specific colors that can also
be used for posters. This one is an ``alert'' block. 
\end{alertblock}
\begin{exampleblock}{Blocks in different colors}

This is an ``example block''.

Yet another color.
\end{exampleblock}
\begin{block}{About logical term markup}

The LyX ``logical markup'' module creates these right-click accessible
types

\structure{structure},\alert{Alert},\code{Code},\noun{Noun},\textbf{Bold},\emph{Emph},\strong{Strong(same as bold)},
and others.

Please be advised that these word-level markers are not to be used
for long text sections. Several words will often not cooperate with
linebreaks. I was worried \strong{Strong(same as bold)} would cause
trouble in that regard.
\end{block}
\end{alertblock}

\column{.49\linewidth}
\begin{block}{Enumerated Lists}
\begin{enumerate}
\item some items 
\begin{enumerate}
\item some items 
\item some items 
\item some items 
\end{enumerate}
\end{enumerate}
%
\begin{block}{Some items and some math}
\begin{itemize}
\item some items and $\alpha=\gamma,\sum_{i}$ 
\item some items 
\item some items 
\end{itemize}
\[
\alpha=\gamma,\sum_{i}
\]

\end{block}
%
\begin{block}{More items \ldots}
\begin{itemize}
\item some items
\item some items
\item some items
\end{itemize}
\end{block}
%
\begin{block}{\ldots{} and yet more!}
\begin{itemize}
\item some items 
\item some items 
\end{itemize}
\end{block}
%
\begin{block}{Contact}
 

\begin{minipage}[t][6cm]{0.2\columnwidth}%
\vspace{-36pt}
\includegraphics[width=5cm]{static/Elmer_Fudd}%
\end{minipage}%
\begin{minipage}[t]{0.8\columnwidth}%
Paul Johnson

University of Kansas

Center for Research Methods and Data Analysis

\url{elmer@wikimedia.org}

Elmer graphic from 

{\footnotesize{}\url{https://commons.wikimedia.org/wiki/File:Elmer_Fudd.png}}{\footnotesize\par}%
\end{minipage}
\end{block}
\end{block}
\end{columns}

\end{frame}
\end{document}
